% authority: science-draft
% status: draft/control
% route: ontology-law-research-packet
% trigger_classification: derivation_critical_missing_source_law
% adoption_status: blocked_adoption_open_continuation
\documentclass[11pt]{article}
\usepackage[margin=1in]{geometry}
\usepackage{amsmath,amssymb,amsthm}
\usepackage{booktabs}
\usepackage{enumitem}

\newtheorem{theorem}{Theorem}
\newtheorem{proposition}{Proposition}
\newtheorem{remark}{Remark}
\newcommand{\Ftwo}{\mathbb F_2}
\newcommand{\Candidate}{\mathsf{EqSrcFlowGeneratedGradedOrbitRootLaw}^{\mathrm{cand},v1}_{\mathrm{src}}}

\title{P1-T01 Smuggling Audit of the Proposal-Only\\
EqSrc Flow-Generated Graded-Orbit Root Law Candidate v1}
\author{Smuggling Auditor control packet RT-20260720-009}
\date{2026-07-20}

\begin{document}
\maketitle

\section{Verdict and authority boundary}

The exact input is the proposal-only candidate artifact with SHA-256
\texttt{b712552d328f144491bff689b702eba6dc2027ce1cc61c7052adbca84b0639f7}.
The audit verdict is
\texttt{conditional\_source\_purity\_pass\_with\_action\_realization\_and\_completeness\_nondiscrimination}.
The displayed reachability, cover, rank-parity, boundary-line, and kernel-pair
claims are mathematically correct after every admitted record field is
supplied. The candidate contains no explicit target-GR import and no explicit
repository-process premise in those definitions.

That conditional pass does not establish source provenance. The ordered
action package can realize every nontrivial ordered finite partition, the
``full component'' condition cannot be certified from the local record alone,
positivity is named but not axiomatized, and the morphism and variation
categories preserve the result by construction. These are source-selection
and robustness obstructions, not evidence of an explicit hidden target import
and not a global no-go theorem.

The candidate remains \texttt{proposal-only} with
\texttt{blocked\_adoption\_open\_continuation}. General EqSrc remains
undischarged. No ontology source, Distance-to-GR ledger, metric-use ledger,
benchmark, or protected decision changes.

\section{Audited proposal data}

The candidate admits a record
\[
 U=(C,\Lambda_+,\oplus,0,\Phi,\mathcal M_U,\mathcal V_U)
\]
only after supplying a monoid action, a full action component, antisymmetric
and locally finite reachability, one global minimum, a common finite
saturated-chain rank, two nonempty parity blocks, and independently declared
morphism and variation categories. It then defines
\[
 P_U=\{x:\rho_U(x)\equiv0\pmod2\},\qquad
 Q_U=\{x:\rho_U(x)\equiv1\pmod2\},
\]
\[
 A_U=\operatorname{span}_{\Ftwo}\{\mathbf1_{P_U},\mathbf1_{Q_U}\},
 \qquad B_U=\Ftwo\mathbf1_{Q_U},
 \qquad aR_Ua'\Longleftrightarrow a-a'\in B_U.
\]
The order of definition is source-facing: no partition or relation is an
explicit input. The audit therefore distinguishes textual source purity from
the stronger question whether the admitted action data select rather than
encode the desired structure.

\section{Universal ordered-partition action realization}

\begin{theorem}[Universal ordered-bipartition action realization]
Let a finite set (C) have a nontrivial ordered partition (C=P\sqcup Q).
There is a finite action record admitted by the candidate whose even and odd
rank blocks are exactly (P) and (Q), respectively.
\end{theorem}

\begin{proof}
Choose (r\in P) and (q_*\in Q). Let the free monoid on symbols
\(\{u_q:q\in Q\}\cup\{v_p:p\in P\setminus\{r\}\}\) act through the
transformations
\[
 u_q(r)=q,\quad u_q(x)=x\ (x\ne r),
 \qquad
 v_p(q_*)=p,\quad v_p(x)=x\ (x\ne q_*).
\]
The resulting reachability order has (r) as its unique minimum. Its covers
are precisely (r\lessdot q) for (q\in Q) and
\(q_*\lessdot p\) for (p\in P\setminus\{r\}). It is finite, locally
finite, antisymmetric, connected under the undirected closure of action
steps, and graded with ranks zero on (r), one on (Q), and two on
\(P\setminus\{r\}). Hence its even block is (P), its odd block is (Q),
and both are nonempty.
\end{proof}

For a fixed ordered partition, the construction has (|P||Q|) choices of
root and bridge. Summed over all ordered nontrivial partitions of an
(n)-element labeled carrier, the audited parameterization count is
\[
 \sum_{\varnothing\ne Q\subsetneq C}|Q|(n-|Q|)
 =n(n-1)2^{n-2}.
\]
For (n=2,3,4,5,6), the partition counts are (2,6,14,30,62), while the
root--bridge parameterization counts are (2,12,48,160,480).

\begin{remark}[Interpretation]
The theorem proves arbitrary source-side encodability under freely selected
action data. It does not prove that the candidate text imports a target
partition, and it does not prove that no conservative source law could select
an action. It shows that the current admission conditions do not discriminate
among the possible ordered partitions.
\end{remark}

\section{Positivity and component-completeness audit}

The symbol \(\Lambda_+\) and the phrase ``source-order monoid'' suggest a
positive sector, but the candidate supplies neither an order relation nor
positive-cone axioms. The action theorems use only a monoid. Thus positivity
is a name rather than a discharged formal condition. This does not invalidate
the conditional algebraic relation; it leaves the claimed source-order
provenance under-specified.

The full-component guard is also not intrinsic to the displayed local tuple.
Consider
\[
 \widetilde C=\{-1,0,1,2,3\},\qquad
 \widetilde\Phi_n(i)=\min(i+n,3).
\]
Its forward-invariant window (C=\{0,1,2,3\}) carries exactly the candidate's
four-event positive action when restricted. Declared in isolation, that
window is internally admitted and rooted at zero. Embedded in
\(\widetilde C\), zero has predecessor (-1), and the full component is
rooted at (-1). No field of the isolated tuple certifies which ambient
completion is authoritative. The candidate therefore fails to distinguish a
full source component from an observationally or procedurally truncated
forward-invariant tail unless external component provenance is supplied.

This pair does not contradict the candidate's conditional theorem: once the
full component is genuinely fixed, its global minimum is natural. It shows
that the admission predicate assumes the very completeness fact needed to
apply the observer-truncation guard.

\section{Morphism, variation, and finite nondiscrimination checks}

The action-naturality theorem is correct for isomorphisms already required to
preserve the action. This is category-by-preservation, not physical covariance.
No independently selected larger source category is tested.
Likewise, a variation is accepted only when it preserves all admission data
up to admitted isomorphism; otherwise it returns bottom. Conditional
fail-closure is coherent, but it is not independent robustness evidence.

On a four-token carrier there are (14) nontrivial ordered partitions and
(24) permutations, giving (336) partition--permutation cases. Exactly
(72) stabilize the chosen odd block and (264) change it. Across all valid
single-token transfers between nonempty blocks, all (48) of (48) change
the boundary line. The candidate's preservation-defined morphisms and
variations exclude rather than resolve these relation-changing branches.

\section{Audit matrix}

\begin{center}
\begin{tabular}{p{0.25\linewidth}p{0.63\linewidth}}
\toprule
Question & Finding \\
\midrule
Ordered-action provenance & Current ontology does not derive the monoid,
action, or action selection. Universal action realization shows admission is
nonselecting across finite ordered partitions. \\
Positivity & No order or positive-cone axiom is defined; the plus sign is
not a proved property. \\
Component selection & ``Full component'' is asserted. A forward-invariant
tail and its ambient completion are locally indistinguishable without extra
provenance. \\
Grading & Rank is well defined conditionally, but the construction can load
any desired parity partition through the selected action. \\
Morphisms & Naturality is valid inside a category defined to preserve the
action; physical covariance is not established. \\
Variations & Preservation or bottom is fail-closed but does not test
independently admissible relation-changing variations. \\
Target or process import & No explicit target-GR, benchmark, validator,
registry, role, handoff, or file-order premise was found in the mathematical
definition. Arbitrary encodability is not proof of an explicit import. \\
\bottomrule
\end{tabular}
\end{center}

\section{Distance-to-GR, freeze review, and next route}

The active burden \texttt{source\_equivalence\_eqsrc} stays at
\texttt{draft object exists}. The audit adds theorem-level obstruction data
but no Distance-to-GR delta. Exact-GR recovery, RetainH, GenH,
\(M_{\rm src}\), \(g_{\rm eff}\), matter coupling, Einstein equations, and
benchmark validation remain open.

Because this is another member of the response-token, orientation-line,
root, partition, and grading selection family, a family-level freeze review
is now required. The exact v21 next work item is \texttt{P1-T02}: one
Ontology Formalizer packet must record stable family identifiers, assumption
hashes, construction--audit--stress--repair lineage, freeze decisions, and
supersession relations. P1-T03, not this audit, decides whether the family is
frozen or one materially distinct theorem route remains.

This packet does not repair or stress the candidate, declare it adopted or
rejected, edit canonical ontology, issue a Gate Chair verdict, or assert
future source-extension impossibility. The candidate remains a
\texttt{draft/control} and \texttt{proposal-only} source-extension object.

\end{document}
